\documentclass{beamer}
\usepackage{etex}

\mode<presentation>
{
  \usetheme{Warsaw}
}


\usepackage[english]{babel}
\usepackage[latin1]{inputenc}
\usepackage{times}
\usepackage[T1]{fontenc}

\usepackage{pstricks}
\usepackage{pst-solides3d}
\usepackage{pstricks-add}
\usepackage{pst-fun}

\usepackage[squaren,Gray]{SIunits}

\usepackage[absolute,overlay]{textpos}

\newfont{\Matrixfont}{necker scaled 2400}
\newfont{\matrixfont}{necker scaled 1200}

\setbeamercolor{framesource}{fg=gray}
\setbeamerfont{framesource}{size=}


\title[SVM for Network AD]{Support Vector Machines for\\ Network Anomaly Detection}
\author{G.~Vormayr}
\institute{University of Technology Vienna}
\date[CNSE 2014]{Communication Networks SE, 2014}
\subject{Support Vector Machines for Network Anomaly Detection}


\pgfdeclareimage[height=0.5cm]{university-logo}{pics/TULogoSW}
\logo{\pgfuseimage{university-logo}}

\newcommand{\mbird}{\mathord{\begin{pspicture}(0,0)(0.7,0.6) \psscalebox{0.2}{\psBird} \end{pspicture}}}
\newcommand{\mfish}{\mathord{\begin{pspicture}(0,0)(0.6,0.4) \psscalebox{0.1}{\psFish[fillstyle=slope]} \end{pspicture}}}


\AtBeginSubsection[]
{
  \begin{frame}<beamer>{Outline}
    \tableofcontents[currentsection,currentsubsection]
  \end{frame}
}

\begin{document}

\begin{frame}
  \titlepage
\end{frame}

\begin{frame}{Outline}
  \tableofcontents
\end{frame}


\section{Support Vector Machines}

\begin{frame}{Motivation}
    \begin{pspicture}(-1,-1)(10,3)
        \psset{lightsrc=-20 0 20,viewpoint=-20 0 3,Decran=50}

        \psset{solidmemory}
        \psset{color1=black!60,color2=black!80}

        %belt
        \pstVerb{/iface 0 store}
        \psSolid[name=top,fcol=20 {iface (color1) iface 1 add (color2) /iface iface 2 add store} repeat,object=grille,ngrid=1 .1,action=none,base=0 0.5 -4 0]
        \pstVerb{/iface1 0 store}
        \psSolid[name=bottom,fcol=20 {iface1 (color2) iface1 1 add (color1) /iface1 iface1 2 add store} repeat,object=grille,ngrid=1 .1,action=none,base=0 0.5 -4 0](0,0,-0.1)

        %rollers
        \psSolid[name=RL,object=cylindre,RotY=90,h=0.5,r=0.05,action=none,fillcolor=black!60](0,0,-0.05)
        \psSolid[name=RR,object=cylindre,RotY=90,h=0.5,r=0.05,action=none,fillcolor=black!60](0,-4,-0.05)

        %detector
        \psSolid[name=detector,object=cube,a=1,action=none,fillcolor=green,incolor=red,hollow,affinage=1 3,chanfrein,chanfreincoeff=.9](0.25,-2,0.2)
        \psSolid[name=display,object=parallelepiped,action=none,fillcolor=black!40,a=0.05,b=0.8,c=0.3,numfaces=all](-0.25,-2,0.45)
        \psSolid[name=knobC,object=tronccone,action=none,r0=0.05,r1=0.04,h=0.1,RotY=-90,fillcolor=yellow](-0.3,-1.8,-0.1)
        \psSolid[name=knobChi,object=tronccone,action=none,r0=0.05,r1=0.04,h=0.1,RotY=-90,fillcolor=red](-0.3,-2.2,-0.1)

        \psSolid[object=fusion,action=draw**,mode=4,base=bottom RL top RR detector display knobC knobChi]
        \composeSolid

        \only<1>{%
            \rput(3.8,0.1){\psscalebox{-0.5 0.5}{\psBird}}
            \rput(1.9,0.1){\psscalebox{-0.12 0.12}{\psFish[fillstyle=slope]}}
            \rput(0.3,0.7){\psscalebox{0.12}{\includegraphics{pics/tux.eps}}}
            \rput(5,1){\textcolor{green}{\Matrixfont{----}}}}
        \only<2>{%
            \rput(3.6,0.1){\psscalebox{-0.12 0.12}{\psFish[fillstyle=slope]}}
            \rput(2.0,0.7){\psscalebox{0.12}{\includegraphics{pics/tux.eps}}}
            \rput(5,1){\textcolor{green}{\Matrixfont{Bird}}}}
        \only<3>{%
            \rput(8.2,0.1){\psscalebox{-0.5 0.5}{\psBird}}
            \rput(3.1,0.7){\psscalebox{0.12}{\includegraphics{pics/tux.eps}}}
            \rput(5,1){\textcolor{green}{\Matrixfont{Fish}}}}
        \only<4>{%
            \rput(9.0,0.1){\psscalebox{-0.5 0.5}{\psBird}}
            \rput(7.1,0.1){\psscalebox{-0.12 0.12}{\psFish[fillstyle=slope]}}
            \rput(5,1){\textcolor{green}{\Matrixfont{????}}}}

        \rput(4.5,0){$C$}
        \rput(5.5,0){$\gamma$}
    \end{pspicture}
    \begin{beamercolorbox}[ht=1cm,right]{framesource}
        \fontsize{3pt}{1em}\selectfont\color{gray} Tux image adapted from: \url{http://www.home.unix-ag.org/simon/penguin/} by Larry Ewing, Simon Budig, Anja Gerwinski, via Wikimedia Commons
    \end{beamercolorbox}
\end{frame}

\begin{frame}
    \begin{block}{What we want}
        \begin{math}
            f: (\mbird, \mfish, \dots) \rightarrow (\text{\matrixfont{BIRD}}, \text{\matrixfont{FISH}}, \dots)
        \end{math}
    \end{block}
    \begin{block}{How we want it}
        \begin{itemize}
            \item Without knowing the relation.
            \item Without learning the relation.
            \item Probably there even is no relation.
        \end{itemize}
    \end{block}
\end{frame}

\subsection{Inner Workings}

\begin{frame}{Features}{Describe the Objects}
    \begin{columns}[t]
        \column[T]{3cm}
            \begin{pspicture}(0,0)(4,3)
                \psscalebox{0.5}{%
                \rput(2.5,2){\psBird}
                \pnode(4.8,1.8){claw}
                \pnode(2.8,2.4){pecker}
                \pnode(4.2,3.2){color}
                \pnode(5,3.1){wing}
                \rput[bl](4,0.2){\rnode{clawl}{\psframebox{claw size}}}
                \rput[bl](0.4,0.8){\rnode{peckerl}{\psframebox{pecker size}}}
                \rput[bl](3,5.4){\rnode{colorl}{\psframebox{color}}}
                \rput[bl](5.5,4.6){\rnode{wingl}{\psframebox{wing size}}}
                \rput[bl](0,3.2){\rnode{sizel}{\psframebox{size}}}
                \rput[bl](6,1.3){\psframebox{weight}}
                \rput[bl](1.2,4.4){\psframebox{speed}}
                \ncline{->}{clawl}{claw}
                \ncline{->}{peckerl}{pecker}
                \ncline{->}{colorl}{color}
                \ncline{->}{wingl}{wing}}
            \end{pspicture}
        \column{6cm}
        \begin{block}{Features}
            \begin{itemize}
                \item Object can be described with features.
                \item Choose distinguishing features.
            \end{itemize}
        \end{block}
    \end{columns}
\end{frame}

\begin{frame}{Feature Vectors}{Convert Features to Numbers}
    \begin{columns}
        \column{.5\textwidth}
            \begin{math}
                v_1 = \left<\ovalnode{VfliesV}{1, 0, 0}, \rnode{VsizeV}{15}, \rnode{VwingsV}{3}, \rnode{VcolorV}{120}\right>
            \end{math}
            \begin{itemize}
                \item \rnode{Vflies}{flies}
                \item $15\centi\meter$ \rnode{Vsize}{size}
                \item $3\centi\meter$ \rnode{Vwings}{wings}
                \item \rnode{Vcolor}{green}
            \end{itemize}
            \nccurve[angleB=-110,nodesep=2pt]{->}{Vflies}{VfliesV}
            \nccurve[angleB=-90,nodesep=2pt]{->}{Vsize}{VsizeV}
            \nccurve[angleB=-90,nodesep=2pt]{->}{Vwings}{VwingsV}
            \nccurve[angleB=-90,nodesep=2pt]{->}{Vcolor}{VcolorV}
        \column{.5\textwidth}
        \begin{block}{Examples}
            \begin{itemize}
                \item Color $\Rightarrow$ Hue
                \item Category $\Rightarrow$ Vector
                    \begin{itemize}
                        \item Flies $\Rightarrow$ $\left<1, 0, 0\right>$
                        \item Swims $\Rightarrow$ $\left<0, 1, 0\right>$
                        \item Crawls $\Rightarrow$ $\left<0, 0, 1\right>$
                    \end{itemize}
            \end{itemize}
        \end{block}
    \end{columns}
\end{frame}

\begin{frame}{Linear Separation}
    \begin{columns}[t]
        \column[T]{4cm}
            \psset{xAxisLabel=a,yAxisLabel=b}
            \begin{psgraph}[axesstyle=frame,subticks=10](0,0)(1,1){4cm}{4cm} 
                \only<1-3>{
                    \psline(-0.1,0.2)(1.1,0.4)
                    \psline[linestyle=dashed](-0.1,0.3)(1.1,0.5)
                    \psline(-0.1,0.4)(1.1,0.6)
                    \psline{<->}(0.8,0.35)(0.767,0.548)
                    \rput[bl](0.79,0.48){\tiny margin}
                }
                \only<1-2>{
                    \fileplot[plotstyle=dots,dotstyle=o]{dataa.dat}
                    \fileplot[plotstyle=dots,dotstyle=triangle]{datab.dat}
                }
                \only<2>{\fileplot[plotstyle=dots,dotstyle=triangle,fillcolor=red]{datac.dat}}
                \only<3>{
                    \fileplot[plotstyle=dots,dotstyle=o,fillcolor=green]{dataa.dat}
                    \fileplot[plotstyle=dots,dotstyle=triangle,fillcolor=blue]{datab.dat}
                }
                \only<4>{
                    \fileplot[plotstyle=dots,dotstyle=triangle]{datab.dat}
                    \psline(-0.1,0.746)(0.746,-0.1)
                    \psline{<->}(0,0)(0.323,0.323)
                    \rput[bl](0.17,0.1){\tiny margin}
                }
            \end{psgraph}
        \column{5cm}
        \begin{block}{\only<1-2>{Training}\only<3>{Usage}\only<4>{One Class}}
            \begin{itemize}
                \only<1-2>{\item<1-2> Maximize margin.
                           \item<2> $C > 0$ allows errors.}
                \only<3>{\item Class can be easily determined.
                         \item Multiple classes with multiple planes/SVMs.}
                \only<4>{\item Maximize margin to origin.
                         \item No labels needed.}
            \end{itemize}
        \end{block}
    \end{columns}
\end{frame}

\begin{frame}{Kernels}{Nonlinear Separation}
    Transform vector to higher dimension.
    \begin{block}{Kernel Functions}
        \begin{itemize}
            \item Linear
            \item Polynomial
            \item Radial Basis Function (RBF)
            \item Sigmoid
        \end{itemize}
    \end{block}
    Those functions need additional parameters!
\end{frame}

\subsection{Step by Step Guide}

\begin{frame}{Feature Selection}
    \begin{itemize}
        \item Don't use too many features.
        \item Select features that distinguish.
        \item Try different feature selections and compare hit rate.
    \end{itemize}
\end{frame}

\begin{frame}{Normalisation}
    \begin{itemize}
        \item Normalize data to range $[-1, 1]$ or $[0, 1]$.
        \item Prevents dominating features.
        \item<2> \alert{Use same method for training and testing!}\\Testing data is allowed to leave the range.
    \end{itemize}
\end{frame}

\begin{frame}{Kernel Selection}
    \begin{itemize}
        \item Start with RBF Kernel.
        \item Only one additional parameter.
        \item May separate linear data better than linear hyperplane.
        \item Try other kernels if testing data matches poorly.
    \end{itemize}
\end{frame}

\begin{frame}{Cross-Validation}{The Better Training}
    \centering
    \begin{pspicture}(8,2)
        \only<1,3>{
            \psframe(0,0)(2,1)
        }
        \only<2>{
            \psframe[fillcolor=red,fillstyle=solid](0,0)(2,1)
        }
        \only<1-2>{
            \psframe(2,0)(4,1)
        }
        \only<3>{
            \psframe[fillcolor=red,fillstyle=solid](2,0)(4,1)
        }
        \psframe(4,0)(6,1)
        \psframe(6,0)(8,1)
        \psbrace[rot=-90,nodesepA=-1cm,nodesepB=-2pt](8,1.05)(0,1.05){Labeled Data}
    \end{pspicture}
    \begin{itemize}
        \item Split labeled data into even sized chunks.
        \item Train with one chunk, test with rest.
        \item Use best results.
    \end{itemize}
\end{frame}

\subsection*{Summary}
\begin{frame}{Summary}
    \begin{itemize}
        \item SVM can classify data without knowing the relation.
        \item<2-| alert@2> Data might be completely unrelated $\Rightarrow$ bogus results.
        \item<3-> Labeled data is needed for training and testing.
        \item<4-> Performance depends on feature selection and parameters.
        \begin{itemize}
            \item<5| alert@5> Don't overdo it!
        \end{itemize}
    \end{itemize}
\end{frame}

\section{Network Anomaly Detection}

\subsection{Challenges}

\begin{frame}
    \begin{gather*}
        v = \left<\text{IP Version}, \text{IHL},\text{DSCP},\text{ECN},\text{Length},\text{Id},\right.\\
        \text{IP Flags},\text{Fragment Offset},\text{TTL},\text{Protocol},\\
        \text{Header Checksum},\text{Source IP Address},\text{Destination IP Address},\\
        \text{Source Port},\text{Destination Port},\text{Sequence Number},\\
        \text{Acknowledgement Number},\text{Data Offset},\text{TCP Flags},\text{Window Size},\\
        \left.\text{TCP Checksum},\text{Urgent Pointer}\right>
    \end{gather*}
    \alert{Good Idea?}
\end{frame}

\begin{frame}{Training Data}
    \begin{itemize}
        \item Labeled data very hard to come by.
        \item Needs constant retraining.
        \item Feature selection hard.
        \item Important features change with time.
    \end{itemize}
\end{frame}


\subsection{Possibilities}

\begin{frame}
    \begin{itemize}
        \item Machine learning approach recognizes unknown anomalies.
        \item Categorize traffic.
        \item Find anomalies in generated data.
        \item Support other methods.
    \end{itemize}
\end{frame}

\begin{frame}
    \begin{itemize}
        \item Support vector machines are fast.
        \begin{itemize}
            \item Use for pre filtering.
            \item Use for aggregation.
        \end{itemize}
        \item Simple to set up.
        \item<2| alert@2> But this can mislead.
    \end{itemize}
\end{frame}


\section*{Summary}

\begin{frame}{Summary}
    \begin{itemize}
        \item Beware of the usage pitfalls.
        \item Use where appropriate.
    \end{itemize}
\end{frame}

\appendix
\section<presentation>*{\appendixname}
\subsection<presentation>*{For Further Reading}

\begin{frame}[allowframebreaks]
    \frametitle<presentation>{For Further Reading}
    
    \begin{thebibliography}{10}
    
    \beamertemplatearticlebibitems
  % Followed by interesting articles. Keep the list short. 
   \bibitem{chang_libsvm:_2011}
       Chang, Chih-Chung and Lin, Chih-Jen
       \newblock LIBSVM: A Library for Support Vector Machines,
       \newblock {\em ACM Trans. Intell. Syst. Technol.}, 27:1-27:27
       2011.

   \bibitem{burges_tutorial_1998}
       Burges, Christopher J. C.
       \newblock A Tutorial on Support Vector Machines for Pattern Recognition
       \newblock {\em Data Mining and Knowledge Discovery}, 121-167
       1998.

  \bibitem{chandola_anomaly_2009}
    Chandola, Banerje, Kumar
    \newblock Anomaly Detection: A Survey
    \newblock {\em ACM Comput. Surv.}, 15:1-15:58,
    2009.

    %TODO
  \end{thebibliography}
\end{frame}

\bgroup

\setbeamercolor{background canvas}{bg=black}
\setbeamertemplate{navigation symbols}{}
\begin{frame}[plain]
\begin{center}
    \includegraphics[width=5cm]{pics/By.eps}\\
    \vspace{2cm}
    \color{white}
    Unless otherwise noted Support Vector Machines for Network Anomaly Detection by
    Gernot Vormayr is licensed under a
    \href{http://creativecommons.org/licenses/by/4.0/}{Creative Commons Attribution 4.0 International License}.
\end{center}
\end{frame}

\egroup

\end{document}

